% Options for packages loaded elsewhere
% Options for packages loaded elsewhere
\PassOptionsToPackage{unicode}{hyperref}
\PassOptionsToPackage{hyphens}{url}
\PassOptionsToPackage{dvipsnames,svgnames,x11names}{xcolor}
%
\documentclass[
  letterpaper,
  DIV=11,
  numbers=noendperiod,
  oneside]{scrartcl}
\usepackage{xcolor}
\usepackage[left=1in,marginparwidth=2.0666666666667in,textwidth=4.1333333333333in,marginparsep=0.3in]{geometry}
\usepackage{amsmath,amssymb}
\setcounter{secnumdepth}{-\maxdimen} % remove section numbering
\usepackage{iftex}
\ifPDFTeX
  \usepackage[T1]{fontenc}
  \usepackage[utf8]{inputenc}
  \usepackage{textcomp} % provide euro and other symbols
\else % if luatex or xetex
  \usepackage{unicode-math} % this also loads fontspec
  \defaultfontfeatures{Scale=MatchLowercase}
  \defaultfontfeatures[\rmfamily]{Ligatures=TeX,Scale=1}
\fi
\usepackage{lmodern}
\ifPDFTeX\else
  % xetex/luatex font selection
\fi
% Use upquote if available, for straight quotes in verbatim environments
\IfFileExists{upquote.sty}{\usepackage{upquote}}{}
\IfFileExists{microtype.sty}{% use microtype if available
  \usepackage[]{microtype}
  \UseMicrotypeSet[protrusion]{basicmath} % disable protrusion for tt fonts
}{}
\makeatletter
\@ifundefined{KOMAClassName}{% if non-KOMA class
  \IfFileExists{parskip.sty}{%
    \usepackage{parskip}
  }{% else
    \setlength{\parindent}{0pt}
    \setlength{\parskip}{6pt plus 2pt minus 1pt}}
}{% if KOMA class
  \KOMAoptions{parskip=half}}
\makeatother
% Make \paragraph and \subparagraph free-standing
\makeatletter
\ifx\paragraph\undefined\else
  \let\oldparagraph\paragraph
  \renewcommand{\paragraph}{
    \@ifstar
      \xxxParagraphStar
      \xxxParagraphNoStar
  }
  \newcommand{\xxxParagraphStar}[1]{\oldparagraph*{#1}\mbox{}}
  \newcommand{\xxxParagraphNoStar}[1]{\oldparagraph{#1}\mbox{}}
\fi
\ifx\subparagraph\undefined\else
  \let\oldsubparagraph\subparagraph
  \renewcommand{\subparagraph}{
    \@ifstar
      \xxxSubParagraphStar
      \xxxSubParagraphNoStar
  }
  \newcommand{\xxxSubParagraphStar}[1]{\oldsubparagraph*{#1}\mbox{}}
  \newcommand{\xxxSubParagraphNoStar}[1]{\oldsubparagraph{#1}\mbox{}}
\fi
\makeatother


\usepackage{longtable,booktabs,array}
\newcounter{none} % for unnumbered tables
\usepackage{calc} % for calculating minipage widths
% Correct order of tables after \paragraph or \subparagraph
\usepackage{etoolbox}
\makeatletter
\patchcmd\longtable{\par}{\if@noskipsec\mbox{}\fi\par}{}{}
\makeatother
% Allow footnotes in longtable head/foot
\IfFileExists{footnotehyper.sty}{\usepackage{footnotehyper}}{\usepackage{footnote}}
\makesavenoteenv{longtable}
\usepackage{graphicx}
\makeatletter
\newsavebox\pandoc@box
\newcommand*\pandocbounded[1]{% scales image to fit in text height/width
  \sbox\pandoc@box{#1}%
  \Gscale@div\@tempa{\textheight}{\dimexpr\ht\pandoc@box+\dp\pandoc@box\relax}%
  \Gscale@div\@tempb{\linewidth}{\wd\pandoc@box}%
  \ifdim\@tempb\p@<\@tempa\p@\let\@tempa\@tempb\fi% select the smaller of both
  \ifdim\@tempa\p@<\p@\scalebox{\@tempa}{\usebox\pandoc@box}%
  \else\usebox{\pandoc@box}%
  \fi%
}
% Set default figure placement to htbp
\def\fps@figure{htbp}
\makeatother





\setlength{\emergencystretch}{3em} % prevent overfull lines

\providecommand{\tightlist}{%
  \setlength{\itemsep}{0pt}\setlength{\parskip}{0pt}}



 


<script>MathJax = { output: { font: 'mathjax-pagella' } };</script>
\KOMAoption{captions}{tableheading}
\usepackage{tikz-cd}
\makeatletter
\@ifpackageloaded{caption}{}{\usepackage{caption}}
\AtBeginDocument{%
\ifdefined\contentsname
  \renewcommand*\contentsname{Table of contents}
\else
  \newcommand\contentsname{Table of contents}
\fi
\ifdefined\listfigurename
  \renewcommand*\listfigurename{List of Figures}
\else
  \newcommand\listfigurename{List of Figures}
\fi
\ifdefined\listtablename
  \renewcommand*\listtablename{List of Tables}
\else
  \newcommand\listtablename{List of Tables}
\fi
\ifdefined\figurename
  \renewcommand*\figurename{Figure}
\else
  \newcommand\figurename{Figure}
\fi
\ifdefined\tablename
  \renewcommand*\tablename{Table}
\else
  \newcommand\tablename{Table}
\fi
}
\@ifpackageloaded{float}{}{\usepackage{float}}
\floatstyle{ruled}
\@ifundefined{c@chapter}{\newfloat{codelisting}{h}{lop}}{\newfloat{codelisting}{h}{lop}[chapter]}
\floatname{codelisting}{Listing}
\newcommand*\listoflistings{\listof{codelisting}{List of Listings}}
\makeatother
\makeatletter
\makeatother
\makeatletter
\@ifpackageloaded{caption}{}{\usepackage{caption}}
\@ifpackageloaded{subcaption}{}{\usepackage{subcaption}}
\makeatother
\makeatletter
\@ifpackageloaded{sidenotes}{}{\usepackage{sidenotes}}
\@ifpackageloaded{marginnote}{}{\usepackage{marginnote}}
\makeatother
\usepackage{bookmark}
\IfFileExists{xurl.sty}{\usepackage{xurl}}{} % add URL line breaks if available
\urlstyle{same}
% fallback for those not using the hyperref driver hyperxmp:
\makeatletter
\@ifundefined{xmpquote}{\newcommand{\xmpquote}[1]{#1}}{}
\makeatother
\hypersetup{
  pdftitle={Universal algebra},
  colorlinks=true,
  linkcolor={blue},
  filecolor={Maroon},
  citecolor={Blue},
  urlcolor={Blue},
  pdfcreator={LaTeX via pandoc}}


\title{Universal algebra}
\usepackage{etoolbox}
\makeatletter
\providecommand{\subtitle}[1]{% add subtitle to \maketitle
  \apptocmd{\@title}{\par {\large #1 \par}}{}{}
}
\makeatother
\subtitle{A group group is an abelian group.}
\author{}
\date{}
\begin{document}
\maketitle


\newcommand{\id}{\text{id}}
\newcommand{\op}{\mathrm{op}}
\newcommand{\Mod}{\mathbf{Mod}}
\newcommand{\C}{\mathcal{C}}

\section{Lawvere theories and their
models}\label{lawvere-theories-and-their-models}

\textbf{Definition.} A \emph{Lawvere theory} is a small category
\(\mathcal T\) with finite coproducts and a specified object
\(T \in \mathcal T\), such that the objects of \(\mathcal T\) precisely
form the set
\(\{0,\, T,\, T + T,\, T + T + T,\, \dots, \sum_{i < n} T, \dots\}\).
For a category \(\mathcal{C}\) with finite products, a
\emph{\(\mathcal{C}\)-model} of a Lawvere theory \(\mathcal T\) is a
functor \(M \colon \mathcal T^\mathrm{op}\to \mathcal{C}\) which
preserves finite products. We denote by
\(\mathbf{Mod}_{\mathcal T}(\mathcal{C})\) the full subcategory of the
functor category \([\mathcal T, \mathcal{C}]\) consisting of all the
\(\mathcal{C}\)-models of \(\mathcal T\). We may drop the mention of
\(\mathcal{C}\) when it is clear from context; we will often take
\(\mathcal{C}= \mathbf{Set}\).

\textbf{Example.} For each natural number \(n \in \mathbb N\), let
\([n] \in \mathbf{Set}\) denote the a set of cardinality \(n\). Consider
the free--forgetful adjunction
\(\mathbf{Set} \substack{\xrightarrow{F} \\ \perp \\ \xleftarrow[U]{}} \mathbf{Grp}\),
and let \(\mathcal T_{\mathrm{Grp}}\) denote the full subcategory of
\(\mathbf{Grp}\) consisting of each \(F[n]\) for all
\(n \in \mathbb N\). Then we have an equivalence of categories
\(\mathbf{Mod}_{\mathcal T_{\mathrm{Grp}}}(\mathbf{Set}) \simeq \mathbf{Grp}\).

First, we observe that \(T_{\mathrm{Grp}}\) is indeed a Lawvere theory.
For each \(n \in \mathbb N\), we have \([n] \cong \sum_{i < n} [1]\) in
\(\mathbf{Set}\). So \(F[n] \cong \sum_{i < n} F[1]\) in
\(\mathbf{Grp}\), since \(F\) preserves coproducts. Thus we get
\(F[n] = \sum_{i < n} F[1]\) in \(\mathcal T_{\mathrm{Grp}}\).

Next, given a \(\mathbf{Set}\)-model
\(M \colon \mathcal T_\mathrm{Grp}^\mathrm{op}\to \mathbf{Set}\), let
\(HM\) be the group with:\\
1. underlying set \(MF[1]\);\\
2. group multiplication
\(M\alpha \colon (MF[1]) \times (MF[1]) \to MF[1]\), where
\(\alpha \colon F[1] \to F[1] + F[1] = F[2]\) is the map in
\(\mathcal T_\mathrm{Grp}\) sending the generator of \(F[1]\) to the
multiplication of the two generators in \(F[2]\);\\
3. identity element \((Me)(*) \in MF[1]\), where \(*\) is the unique
element in \((MF[1])^0\), and \(e \colon F[1] \to F[0]\) is the unique
map in \(\mathcal T_\mathrm{Grp}\) from \(F[1]\) to \(F[0]\).\\
We extend \(H\) to a functor
\(H \colon \mathbf{Mod}_{\mathcal T_{\mathrm{Grp}}}(\mathbf{Set}) \to \mathbf{Grp}\)
as follows: given another
\(M' \in \mathbf{Mod}_{\mathcal T_\mathrm{Grp}}(\mathbf{Set})\) and a
morphism \(\mu \colon M \to M'\) in
\(\mathbf{Mod}_{\mathcal T_\mathrm{Grp}}(\mathbf{Set})\), we let
\(H\mu\) be the function \(\mu_{F[1]} \colon MF[1] \to M'F[1]\). \[
\begin{tikzcd}
    {MF[2]} && {MF[1]} \\
    \\
    {M'F[2]} && {M'F[1]}
    \arrow["{M\alpha}", from=1-1, to=1-3]
    \arrow["{\beta_{F[2]}}"', from=1-1, to=3-1]
    \arrow["{\beta_{F[1]}}", from=1-3, to=3-3]
    \arrow["{M'\alpha}"', from=3-1, to=3-3]
\end{tikzcd}
\]

Finally, given a group \(G\), define a functor
\(KG \colon \mathcal T_\mathrm{Grp}^\mathrm{op}\to \mathbf{Set}\) by:\\
1. \((KG)(F[n]) xx\)

Details.




\end{document}
